% §4 Scene Belief Reconstruction
% Target: ~1.5 pages for NeurIPS Datasets & Benchmarks track

\section{Scene Belief Reconstruction}
\label{sec:reconstruction}

Question-level accuracy tells us how often a model's answers match the ground truth, but not whether those answers are mutually consistent or whether they imply a coherent spatial layout. To move beyond per-question scoring, the benchmark provides a \emph{scene belief reconstruction} pipeline that converts a model's discrete responses into spatial constraints and solves for the 2D scene layout that best satisfies them. The resulting reconstruction---or its failure---reveals whether a model's answers can be assembled into a geometrically plausible world.

\subsection{From Responses to Constraints}
\label{sec:constraints}

Each model response to a QRR or TRR probe is converted into a spatial constraint over object positions.

\paragraph{Distance poset $\mathcal{P}_{\mathrm{dist}}$.}
Each QRR response contributes an ordinal constraint on pairwise distances. For a probe comparing pairs $(A,B)$ and $(C,D)$:
\begin{itemize}[nosep]
  \item Response ``$<$'' $\;\Rightarrow\;$ $d(A,B) < d(C,D)$
  \item Response ``$\approx$'' $\;\Rightarrow\;$ $d(A,B) \approx d(C,D)$
  \item Response ``$>$'' $\;\Rightarrow\;$ $d(A,B) > d(C,D)$
\end{itemize}
Pairs declared approximately equal are merged into equivalence classes via union-find, and strict orderings form directed edges. The resulting structure is a partially ordered set (poset) on distance equivalence classes, organized as a directed acyclic graph (DAG). A cycle in this DAG indicates a logical inconsistency: the model's distance comparisons are mutually contradictory.

\paragraph{Angular sector system $\mathcal{P}_{\mathrm{ang}}$.}
Each TRR response specifies a clock-hour direction for the target relative to a reference frame defined by two other objects. This maps to an angular arc interval: an hour-level response ($30°$ sector) constrains the relative angle $\alpha$ (Eq.~\ref{eq:trr-angle}) to lie within $\pm 15°$ of the expected direction. When multiple TRR responses constrain the same object triple, the feasible region is the intersection of their arcs; an empty intersection indicates conflicting directional judgments.

\paragraph{Constraint extraction modes.}
The benchmark supports two extraction modes. In the \emph{ground-truth} mode, only correctly answered probes contribute constraints, testing whether accurate sub-answers suffice for global reconstruction. In the \emph{belief} mode, all model predictions (correct or not) are used, reconstructing the model's \emph{externalized spatial belief} regardless of correctness---analogous to how behavioral experiments infer latent representations from response patterns rather than accuracy alone.

\subsection{Optimization}
\label{sec:optimization}

Given the constraint sets $\mathcal{P}_{\mathrm{dist}}$ and $\mathcal{P}_{\mathrm{ang}}$, the benchmark solves for 2D object positions $\{(x_i, y_i)\}_{i=1}^{n}$ that best satisfy all constraints. The optimization uses L-BFGS-B with a gauge-fixing scheme that removes the inherent ambiguities of ordinal data.

\paragraph{Gauge fixing.}
Ordinal constraints determine positions only up to a similarity transformation (translation, rotation, scale) and reflection. To resolve these five degrees of freedom, three \emph{anchor} objects are selected by constraint participation frequency and fixed:
\begin{equation}
  \mathbf{x}_a = (0, 0), \quad \mathbf{x}_b = (1, 0), \quad y_c \geq 0.
  \label{eq:gauge}
\end{equation}
The first two anchors fix translation, rotation, and scale (4~DOF); the inequality on $y_c$ breaks the reflection ambiguity (1~DOF). The remaining $2(n{-}2)$ coordinates are free variables optimized via L-BFGS-B, with $y_c \geq 0$ enforced as a box constraint.

\paragraph{Loss function.}
The total objective combines three terms: $\mathcal{L} = \mathcal{L}_{\mathrm{QRR}} + \mathcal{L}_{\mathrm{TRR}} + \mathcal{L}_{\mathrm{sep}}$.

For QRR constraints, we use a log-domain ranking loss. Let $\delta = \log d(A,B) - \log d(C,D)$. For a ``$<$'' constraint, the loss is a softplus penalty when the margin is violated:
\begin{equation}
  \mathcal{L}_{\mathrm{QRR}}^{(<)} = \frac{1}{\beta}\log\bigl(1 + \exp\bigl(\beta(\delta + m)\bigr)\bigr),
  \label{eq:qrr-loss}
\end{equation}
where $m = 0.1$ is the margin and $\beta = 10$ controls sharpness. The ``$>$'' case is symmetric ($-\delta + m$), and ``$\approx$'' uses a Huber loss on $\delta$ with tolerance $\delta_{\mathrm{eq}} = 0.1$. The log domain ensures scale invariance: the loss depends only on distance \emph{ratios}, matching the ordinal nature of QRR comparisons.

For TRR constraints, we use a sector-tolerance loss. Let $\mathbf{u}$ be the unit vector from \texttt{ref1} to \texttt{ref2} (the 12~o'clock direction), $\mathbf{v}$ the unit vector from \texttt{ref1} to the target, and $\mathbf{u}_\alpha = R(\alpha)\mathbf{u}$ the expected direction for clock hour $h$ (where $\alpha = 30° \cdot (h \bmod 12)$). The loss penalizes configurations where the cosine similarity falls below the sector boundary:
\begin{equation}
  \mathcal{L}_{\mathrm{TRR}} = \frac{1}{\beta}\log\bigl(1 + \exp\bigl(\beta \cdot \tfrac{\cos\theta_{\mathrm{tol}} - \mathbf{u}_\alpha \cdot \mathbf{v}}{\tau_s}\bigr)\bigr),
  \label{eq:trr-loss}
\end{equation}
where $\theta_{\mathrm{tol}} = 15°$ for hour-level and $45°$ for quadrant-level constraints, and $\tau_s = 0.1$ is a softening temperature (distinct from the QRR comparator tolerance $\tau$ in \S\ref{sec:tasks}).

A separation regularizer $\mathcal{L}_{\mathrm{sep}} = \lambda \sum_{i < j} \max(0, \epsilon - \|\mathbf{x}_i - \mathbf{x}_j\|)$ with $\epsilon = 0.05$ and $\lambda = 1.0$ prevents object collapse.

\paragraph{Multi-start optimization.}
Because the loss landscape may contain local minima (especially for complex scenes or inconsistent constraints), the solver runs $K = 10$ restarts from random initializations and collects all converged solutions, sorted by loss. The diversity of solutions is itself informative: a unique minimum suggests the constraints tightly determine the layout, while multiple distinct minima indicate geometric ambiguity.

\subsection{Scene-Level Metrics}
\label{sec:recon-metrics}

The reconstruction produces a suite of scene-level metrics that go beyond per-question accuracy.

\paragraph{Constraint Satisfaction Rate (CSR).}
The fraction of input constraints satisfied by the best reconstruction. CSR$_{\mathrm{QRR}}$ checks whether the reconstructed pairwise distances preserve each ordinal comparison (using the same $\tau = 0.10$ tolerance); CSR$_{\mathrm{TRR}}$ checks whether reconstructed angles fall within the specified sector. A CSR near 1.0 means the model's answers are mutually consistent; a low CSR means no single layout can reconcile them.

\paragraph{Normalized RMS Error (NRMS).}
When ground-truth positions are available, the best reconstruction is Procrustes-aligned (translation, rotation, uniform scale) to the ground truth, and the root-mean-square positional error is normalized by the spatial extent of the ground-truth layout:
\begin{equation}
  \mathrm{NRMS} = \frac{\mathrm{RMS}(\hat{\mathbf{X}}_{\mathrm{aligned}}, \mathbf{X}_{\mathrm{GT}})}{\max_k \operatorname{range}_k(\mathbf{X}_{\mathrm{GT}})},
  \label{eq:nrms}
\end{equation}
where $\operatorname{range}_k$ denotes the extent along coordinate axis $k$. NRMS $\approx 0$ indicates a faithful reconstruction; NRMS $> 0.3$ typically indicates severe geometric distortion.

\paragraph{Kendall $\tau$.}
The Kendall rank correlation between reconstructed and ground-truth pairwise distances measures how well the ordinal distance structure is preserved, independent of metric accuracy. A $\tau$ near 1.0 means the model preserves the ``which pair is closer'' ordering; a $\tau$ near 0 indicates no ordinal structure.

\paragraph{Geometric multiplicity ($K_{\mathrm{geom}}$) and spread.}
From the multi-start solutions, we cluster converged results by RMS distance (threshold 0.10) after filtering solutions with loss exceeding $3\times$ the best. $K_{\mathrm{geom}}$ is the number of distinct clusters (geometric modes). A \emph{spread} metric quantifies the average RMS distance of good solutions from the best, measuring how tightly the constraints pin down the layout. $K_{\mathrm{geom}} = 1$ with low spread indicates a well-determined reconstruction; $K_{\mathrm{geom}} > 1$ indicates the constraints admit multiple qualitatively different layouts.

\subsection{Validation}
\label{sec:validation}

To verify the reconstruction pipeline, it is run on ground-truth constraints extracted directly from the scene annotations (i.e., all QRR and TRR answers are correct). Under these ideal inputs, the solver consistently recovers layouts with CSR$_{\mathrm{QRR}} \approx 1.0$, NRMS $\approx 0.01$, and $K_{\mathrm{geom}} = 1$ across scenes of all complexity levels ($n = 4$ to $n = 10$). This confirms that the loss functions and gauge-fixing scheme faithfully recover the ground-truth geometry when the input constraints are consistent, establishing a sound baseline for interpreting VLM reconstruction results: any degradation in reconstruction quality can be attributed to errors or inconsistencies in the model's responses, not to limitations of the reconstruction method itself.
